\documentclass[conference, 10pt]{IEEEtran}
\IEEEoverridecommandlockouts
% The preceding line is only needed to identify funding in the first footnote. If that is unneeded, please comment it out.




%%%%%
%% Packages:
\usepackage[utf8]{inputenc} 
\usepackage[T1]{fontenc}
\usepackage{url}              % provides \url{...}
%\usepackage{ifthen}          % provides \ifthenelse
\usepackage{cite}             % improves presentation of citations

\usepackage[cmex10]{amsmath}  % Use the [cmex10] option to ensure complicance
                              % with IEEEXplore (see bare_conf.tex)
\interdisplaylinepenalty=1000 % As explained in bare_conf.tex
\usepackage{mleftright}       % fix to wrong spacing of \left-,
\mleftright                   % \middle- \right-commands 
\usepackage{amsthm}
\usepackage{graphicx}         % provides \includegraphics{...} to
                              % include graphics (pdf format)
\usepackage{booktabs}         % fixes poor spacing in tables and
                              % provides \toprule, \midrule, \bottomrule

%\usepackage{algorithmicx}    % provides an algorithmic environment for
                              % describing algorithms. See
                              % https://ctan.org/pkg/algorithmicx

% \usepackage[caption=false,font=footnotesize]{subfig}
                              % provides subnumbering within a
                              % floating figure or table

%% For arrays and multiple-line equations, use the
%% IEEEeqnarray-environment. See
%%              https://moser-isi.ethz.ch/manuals.html#eqlatex  
%% for instructions.

%% Do NOT use amsthm or hyperref!
%% -IEEEtran provides its own versions of theorems.
%% -IEEEXplore does not accept submissions with hyperlinks


%%%%%
%% correct bad hyphenation here
\hyphenation{op-tical net-works semi-conduc-tor}

% -------------------------------------------------------------------------


%%%%%%%%%%%%%%%%%%%%%%%%%%%%%%%%%%%%%
%% Packages and commands used in this work 
%%%%%%%%%%%%%%%%%%%%%%%%%%%%%%%%%%%%%
\usepackage{cite}
\usepackage{amsmath}
\usepackage{amssymb}
\usepackage{algorithmic}
\usepackage{textcomp}
\usepackage{xcolor}
\def\BibTeX{{\rm B\kern-.05em{\sc i\kern-.025em b}\kern-.08em
    T\kern-.1667em\lower.7ex\hbox{E}\kern-.125emX}}
% \usepackage{subcaption}
\usepackage{graphicx}
\usepackage{subfiles}
\usepackage[inline]{enumitem}
\newcommand{\bvec}[1]{\ensuremath{\mathbf{#1}}}
%\usepackage{arydshln} %for dashed line in  atable
\usepackage{array} 
\usepackage{url}
\usepackage{epstopdf}
\usepackage{tikz}
\usetikzlibrary{automata,arrows,positioning,calc,
shapes,chains}
\usepackage{mathtools}
\usepackage{texdef2015}
\allowdisplaybreaks
\newcommand{\nstack}[2]{\begin{matrix} {\mathbf #1}\\ (#2)\end{matrix}}
\newcommand{\au}{a^{(u)}}
\newcommand{\AUcal}{{\Acal}^{(u)}}
\newcommand{\Lcalhat}{\hat{\mathcal{L}}}
\newcommand{\LUcal}{{\Lcal}^{+}}
\newcommand{\Lcal}{\mathcal{L}}
\newcommand{\Qcal}{\mathcal{Q}}
\newcommand{\Acal}{\mathcal{A}}
\newcommand{\Bcal}{\mathcal{B}}
\newcommand{\R}{\mathbb{R}}
\newcommand{\dq}[2][q]{\delta_{#2,#1}}
\newcommand{\dqt}[2][q(t)]{\delta_{#2,#1}}
\newcommand{\ql}{q_l}
\newcommand{\qpl}{q'_l}
\newcommand{\laml}[1][l]{\lambda^{(#1)}}
\newcommand{\muinv}{\mu^{-1}}
\newcommand{\psim}[2][m]{\psi^{(#1)}_{#2}}
\newcommand{\psis}[2][s]{\psi^{(#1)}_{#2}}
\newcommand{\Lpsis}[2][s]{L\psis[#1]{#2}}
\newcommand{\Lpsim}[2][m]{L\psim[#1]{#2}}
\newcommand{\Lamm}[2][j]{\Lambda_{#2#1}}
\newcommand{\cvec}[1]{\Bracket{#1}}
\newcommand{\cvect}[1]{\cvec{#1}^\top}
\newcommand{\smvec}[1]{\left[\begin{smallmatrix} #1\end{smallmatrix}\right]}
\newcommand{\rowvec}[1]{[\begin{matrix} #1\end{matrix}]}
\newcommand{\rvec}[1]{\smrvec{#1}}
\newcommand{\smrvec}[1]{\setlength\arraycolsep{2pt}\rowvec{#1}}
\newcommand{\avgage}{\tilde{\age}}
\newcommand{\piv}{\text{\boldmath{$\pi$}}}
\newcommand{\pivbar}{\text{\boldmath{$\bar{\pi}$}}}
\newcommand{\pibar}{\bar{\pi}}
\newcommand{\vvbar}{\bar{\vv}}
\newcommand{\vvhat}{\hat{\vv}}
\newcommand{\uvhat}{\hat{\uv}}
\newcommand{\agesum}{\age_{\Sigma}}
\newcommand{\zerov}[1][\mbox{}]{\mathbf{0}_{#1}}
\newcommand{\onev}[1][n]{{\mathbf{1}}_{#1}}
\renewcommand{\vec}[1]{\begin{bmatrix} #1\end{bmatrix}}
\renewcommand{\bv}{\mathbf{b}}%Texdef makes this vector v !
\newcommand{\range}[2][0]{#1\,{:}\,#2}
\newcommand{\vm}[2][m]{v^{(#1)}_{#2}}
\newcommand{\vs}[2][s]{v^{(#1)}_{#2}}
\newcommand{\vmdot}[2][m]{\dot{v}^{(#1)}_{#2}}
\newcommand{\vsdot}[2][s]{\dot{v}^{(#1)}_{#2}}
\newcommand{\vmbar}[2][m]{\bar{v}^{(#1)}_{#2}}
\newcommand{\vsbar}[2][s]{\bar{v}^{(#1)}_{#2}}
\newcommand{\vvm}[2][m]{\vv^{(#1)}_{#2}}
\newcommand{\vvs}[2][s]{\vv^{(#1)}_{#2}}
\newcommand{\vvmdot}[2][m]{\dot{\vv}^{(#1)}_{#2}}
\newcommand{\vvsdot}[2][s]{\dot{\vv}^{(#1)}_{#2}}
\newcommand{\vvmbar}[2][m]{\bar{\vv}_{#2}}
\newcommand{\vvmhat}[2][m]{\hat{\vv}_{#2}}
\newcommand{\vvsbar}[2][s]{\bar{\vv}^{(#1)}_{#2}}
\newcommand{\first}{1}
\newcommand{\nlen}{n} 
\newcommand{\smeval}[1]{#1|}
\newcommand{\Tcal}{\mathcal{T}}
\newcommand{\age}{\Delta}
\newcommand{\Spre}{Z}
\newlength{\swwidth}
\newcommand{\stw}[2]{\text{\settowidth{\swwidth}{\text{\ensuremath{#1}}}\makebox[\swwidth]{#2}}}
\newcommand{\vz}[1][0]{\stw{v_{00}}{#1}}
\newcommand{\xz}[1][0]{\stw{x_{0}}{#1}}
\newcommand{\Yq}{Y^{(q)}}
\newcommand{\DotM}{$\cdot$/M}
\newcommand{\Pc}{P_c}
\newcommand{\Pe}{P_e}
\newcommand{\slotage}{\age_{\text{slotted}}}

\newcommand{\al}{\alpha}
\newcommand{\alhat}{\hat{\alpha}}
\newcommand{\lam}{\lambda}
\newcommand{\lamhat}{\hat{\lambda}}
\newcommand{\lamsum}{\lambda^*}
\newcommand{\om}{\omega}
\newcommand{\sig}{\sigma}
\newcommand{\gam}{\gamma}
\newcommand{\mh}{\hat{\mu}}
\newcommand{\musum}{\mu^*}
\newcommand{\bt}{\beta}
\newcommand{\xh}[1]{\hat{x}_{#1}}
\newcommand{\rh}{\hat{\rho}}

\newcommand{\oneage}{x_1}
\newcommand{\twoage}{x_2}
\newcommand{\maxage}{\hat{x}}
\newcommand{\dpage}{\hat{y}}
\newcommand{\monage}{\hat{z}}

\newcommand{\mathitem}{\hspace*{1.2em}&\bullet\hspace*{\labelsep}}

\newtheorem{corollary}{Corollary}[theorem] % Use theorem counter as `parent`

\newcommand{\RY}[1]{\quad {\color{red} #1} \quad}
\newcommand{\V}[1]{\quad{\color{blue} #1}\quad}
\newcommand{\MOVE}[2][\mbox{}]{{\color{purple}MOVE: {\color{blue} #1}\ #2}\newline}
\setlength{\footskip}{12.3pt}
\newtheorem*{remark}{Remark}

\begin{document}

\title{Timely Processing Of Updates From Multiple Sources} 

% %%%%%%
% \author{%
%   \IEEEauthorblockN{Anonymous Authors}
%   \IEEEauthorblockA{%
%     Please do NOT provide authors' names and affiliations\\
%     in the paper submitted for review, but keep this placeholder.\\
%     ISIT23 follows a \textbf{double-blind reviewing policy}.}
% }

\author{\IEEEauthorblockN{Vishakha Ramani,
Ivan Seskar,
Roy D. Yates}
\IEEEauthorblockA{WINLAB, Rutgers University \\
Email: \{vishakha, seskar, ryates\}@winlab.rutgers.edu
}}


\maketitle

\begin{abstract}
We consider a system where the updates from independent sources are disseminated via a 
% shared memory publication mechanism 
publish-subscribe mechanism. The sources are the publishers
and a decision process (DP), acting as a subscriber, derives decision updates from the source data.
We derive the stationary expected age of information (AoI) of decision updates delivered to a monitor. 
We show that a lazy computation policy in which the DP may sit idle before computing its next decision update can reduce the average AoI at the monitor even though the DP exerts no control over the generation of source updates. 
This AoI reduction is shown to occur because lazy computation can offset the negative effect of high variance in the computation time.  
\end{abstract}

\section{Introduction}
A dense metropolis is a complex traffic environment, and autonomous cars with a plethora of attached sensors still have limited situational awareness. Hence, holistic situational awareness of a cloud-connected vehicle is facilitated by
\begin{enumerate}
    \item timely collection of  sensory inputs from different sources (e.g., other vehicles, pedestrians, and smart city infrastructure sensors)
    % in a shared database, 
    , and
    \item employing this sensor data 
    % for a potential collision
    to provide timely feedback or decision updates to participating mobile clients. 
\end{enumerate}
For example, cameras at a smart-city intersection \cite{Cosmos} can capture video or images of the intersection and send this data to a processing system at the edge. The system can then analyze the data to identify and inform potential public safety risks such as vehicles running red lights or pedestrians crossing the street outside of designated areas.
% For example,  
% an edge cloud can construct and store 3D feature maps
% from offloaded images 
% of participating clients.
% This can be used to provide accurate localization services to clients passing through the region. The edge cloud can also run object detection services to detect locations of  new objects in the image frame which are then 
% shared with nearby clients~\cite{luyang-edgesharing}.
% These locations can then be shared with nearby clients to extend their awareness \cite{luyang-edgesharing}.

Although, it is possible that the system delivers source updates to interested clients using traditional synchronous request/reply communication paradigm, the rigid structure of such interaction renders the system inefficient for real-time decisions \cite{Carzaniga-pubsub}.
% In such edge computing applications, the implementation of 
Publish-Subscribe (Pub-Sub) systems are an 
alternative communication paradigm 
% common approach 
that enables efficient and scalable communication among various components. The delivery of information from publishers to subscribers is decoupled, meaning that they need not  be aware of each other's existence.
For example, cameras can publish video frames to the system while a processing system subscribes to this data and uses it to make real-time decisions about intersection safety, such as adjusting the timing of traffic lights or alerting drivers to pedestrians in the area.
% Instead, 
A middleware (also known as a broker) acts as an intermediary 
%VR: The next few words can be removed??, using a real-time database 
to manage this distribution of information from publishers to subscribers.

The presence of different types of subscribers with varying time scales of operation in an edge computing system allows for flexibility and adaptability to different requirements and use cases. 
% Different types of subscribers may operate on different time scales based on their requirements and use cases. 
For instance, a real-time subscriber could be an analytics module that processes sensor data in real-time to detect anomalies or trigger immediate actions. On the other hand, a batch processing subscriber could be a processing module that handles a batch of sensor data collected over a certain time window to generate periodic reports or perform long-term trend analysis.
Consequently, such heterogeneity requires decoupling subscribers from publishers. 
% Publish-subscribe messaging can be used to efficiently transfer data between cameras and the processing system. 

In order to facilitate the required information dissemination,
the data is pushed to the middleware from publisher. However, 
two complementary communication modes between middleware and subscriber exist: push and pull modes. Under the push mode, the subscriber passively receives information from the pub-sub middleware. In a pull-based approach, subscribers request messages from the middleware when they are ready to receive them, rather than waiting for the middleware to push messages to them.

There are pros and cons to both approaches. However, the pull mode is better suited for handling a diverse range of subscribers, without requiring a broker to determine the data transfer rate for each of them.
Subscribers have more control over the rate at which they consume messages, making it easier to manage their individual needs \cite{kafka}. 
% On the contrary, the push mode sometimes is cumbersome because of different 

% \V{On pull based vs push based: Some people have criticised this paper that we're still working on because of this asynchronous pull-based subscriber model. What is a favorable design choice?
% Disadvanatges of push-based: a lot of state that needs to be saved, desired schedule, mode of operation, frequency
% In the pull-based, the system should work fine if the subscriber checks often enough. }

Along with this, a real-time database system, acting as a middleware system,
can support information dissemination between publishers
% , such as IoT devices or sensors, 
and subscribers.
% , such as edge servers or analytics modules.
With these design choices in mind, we present an analytical framework 
% for efficient publish-subscribe communication paradigm for disseminating the status information of status updates.
% This work aims to develop an analytical framework, 
based on the publish-subscribe interaction paradigm, that is to be used for disseminating source updates to interested clients that process these source updates and arrive at decisions in a timely way.
% For example, if the received video frame is substantially outdated, then subscribers using obsolete frame might make inappropriate decisions.
% Pub-sub systems in edge computing can utilize memory-based communication mechanisms, such as distributed shared memory or shared buffers, to facilitate efficient message/event exchange between publishers and subscribers. Edge devices, such as IoT devices or sensors, can act as publishers, generating updates. Subscribers, such as edge servers or analytics modules, can subscribe to these updates in order to process and analyze the data in real-time.

\textbf{Considered Problem: } 
The freshness of status information received by subscriber plays an important role in decision making. 
In this work, we focus on a  fundamental problem: What is the average Age of Information (AoI) \cite{Kaul-YG-infocom2012} of  decision updates that are computed from  time-varying set of sensor data published in the database\footnote{Going forward, we refer to the database as {\em shared memory}, or simply as {\em memory}.}?
We model a  class of systems (see Fig.~\ref{fig:tc-sys-model}) in which two independent sources submit time-stamped updates to a  writer that is responsible for publishing the source measurements as updates in the memory.
A decision process (DP), as a subscriber, reads the pair of source $1$ and source $2$ updates from memory and derives a computational result, a {\em decision update}, from this pair that is delivered to a monitor. 

\begin{figure}[t]
    \centering
    \includegraphics[width = \linewidth]{Figures/tc-sys-model.png}    \caption{A writer updates shared database with information fetched from two external sources. A decision process (DP) requests a reader process to read the pair of source updates from the memory. Monitors that track the age of source $1$ and $2$ updates  in the memory are denoted $\bullet(1)$ and $\bullet(2)$ respectively; $\bullet(\maxage(t))$ tracks the age of max-age process in the memory, $\bullet(\dpage(t))$ tracks the age of sampled max-age process, and $\bullet(\monage(t))$ tracks the age of computed decision updates at the external monitor.}
    \label{fig:tc-sys-model}
\end{figure} 


\textbf{Related Work: }
% Periodic transactions updating real time databases in \cite{liu-multiversion, ramamritham-deadlines} was perhaps the earliest introduction of the concepts of age and dispersion of data accessed by the real-time transactions. 
% Authors in \cite{liu-multiversion} performed a simulation based study on the effect of concurrency constructs applied to shared memory and their effects on the age of source and derived data. 
The issue of update timeliness using Age of Information (AoI) metric has been extensively studied; see the surveys  \cite{kosta2017age,Yates-SBKMU-2021jsac-survey} and references therein. A majority of studies focus on analyzing the average AoI of different sources in single-server multi-source queueing models under different service policies \cite{moltafet-multi-source, Yates-Kaul-isit2012, Huang-Modiano-isit2015}. 

More recently, there have been efforts to employ age optimization in diverse applications such as edge cloud processing for low-latency edge-assisted applications \cite{Zhong-ZYGZ-aoi2019, Yates-THR-infocom2017}, and timely mobile routing \cite{ramani-cyINFOCOM, ramani-cy-AoI-INFOCOM}. 
Authors in \cite{Zhong-ZYGZ-aoi2019} discuss a greedy traffic scheduling policy that selects the next processing request (job) that offers maximum age penalty reduction. \cite{Yates-THR-infocom2017} provides an analytical framework for the problem of optimizing frame rate and lag synchronization of server and player in a real-time cloud-assisted gaming application. 
Authors in \cite{ramani-cyINFOCOM, ramani-cy-AoI-INFOCOM} study the effects of concurrency constructs on timely updating of shared data structures used in network software and how this, in turn, affects timely routing of information updates.
% Authors in \cite{Alabbasi2022} study scheduling algorithms to jointly optimize information freshness and job completion time where multiple IoT devices update information on the cloud, and request a computation from the cloud at certain times. 


In previous works, the system models assumed a tight-coupling between the source and monitor. The primary goal was to maintain the status information of the physical process at the destination nodes. For instance, the just-in-time update policy \cite{Kaul-YG-infocom2012} generated a fresh update instantaneously by the source, which started serving as soon as the current update in service was delivered to the destination node. Many works \cite{Zhou-Saad-arxiv2018, Bedewy-SS-arxiv2017old, Yates-isit2015, Sun-UBYKS-IT2017UpdateorWait} have extended upon this idea and aimed to find optimal policies for determining sampling times and updating processes to minimize the Age of Information (AoI) at the destination node.

% For the systems where the producers of source updates and 
In publish-subscribe systems with decoupled producers and consumers, where producers use independent status updating policies, and consumers compute on this status information, a key open problem is to study optimal policies for minimizing the Age of Information (AoI) of computed updates. This work aims to study how the decoupling of publishers and subscribers affect the
timeliness of updates computed 
from base set of sensor (source) updates.
% \cite{Bastopcu-Tnetwork2021, Ulukus-ITW2019, Buyukates-Tcomm2020}. 

% Additionally, concurrency constructs are an important design choice in shared memory systems. 



\section{System Overview}
\label{sec:sysmodel}

% This shared memory model might not be the best design choice for a single reader-writer system, as the communication between the writer and reader processes can be facilitated by using Single-Producer/Single-Consumer queue \cite{Lamport-Reading-Writing}. However, in systems that need to scale to large numbers of subscribers (readers) to source data, using shared memory to provide concurrent access to various reader processes has been recognized as a good design choice \cite{Attiya-smmpi-ibm}. 
% Additionally, this shared memory model provides an abstraction for distributed shared memory systems where the writers and readers are not necessarily co-located. The advantage of using such a paradigm in a practical system is that it provides greater network scalability and flexibility in terms of network topology. 
% In other words, the writers and readers are spatially and temporally loosely coupled.
In this work, we employ a model that captures  the asynchronous operation of the writer and reader of updates in the shared memory Pub-Sub system. 
There are three aspects to the system depicted in Fig.~\ref{fig:tc-sys-model}: 
\begin{enumerate*}
    \item writing the time-varying data received from two sources into the memory,
    \item the arbitration between reader and writer to access memory,
    \item reading the source data from memory and generating a decision update.  
\end{enumerate*}
We now give a brief overview of the writing, reading and decision computation processes.
%, followed by our choice of concurrency construct. 
% We first describe our choice of concurrency construct, followed by modelling the writing and reading processes to analyze the system in terms of AoI freshness metric.

\subsection{Writing source updates to the memory}
We assume each source $i \in \set{1,2}$ independently submits updates as a rate $\lam_i$ Poisson process to the network and 
%We assume 
that these updates arrive fresh at the writer, i.e. with age $0$. 
The write operations to memory have independent exponential $(\mu)$ service times.
We model the writer as a buffer-less service facility with blocking discipline.  Under this model, a source update arriving at the writer will be served only if the writer is idle; otherwise, the update is discarded.

\textit{Remark 1: } In the present study, we investigate a computational regime characterized by relatively longer decision update times compared to the write times of any update in the memory. Our focus is not on regimes where writing to the memory is the overloaded process. Instead, we are primarily interested in examining the delays associated with computational processing. Whether we adopt a buffer-less or a queuing model, the impact of queuing at the writer is expected to be minimal.

%RY This should go somewhere:
%The aforementioned system characteristics dictate that the data held by decision process might be stale compared to data in the memory. 

\subsection{Read-Copy-Update}
We assume that arbitration between the reader and writer for the shared memory is facilitated by Read-Copy-Update (RCU) mechanism. RCU is a lock-less concurrency construct that allows concurrent forward progress for both reader and writer. \cite{MckenneyRCU2001}. RCU can be broadly described in two steps \cite{kernelRCU}: 
\begin{enumerate*}
     \item To publish a newer version of a data item,
     % \footnote{A writer has \textit{published} an update when it exits its write critical section}
     the writer creates a copy of the RCU protected data item, modifies this copy with the newer version of this data item, and atomically replaces the old reference with a reference to this newer version.
     %of data item.
     This publishing process runs concurrently 
     % \RY{and asynchronously?} 
     with ongoing read processes that continue to read the old copy/version using the old reference. However, new read requests 
     %after the publication 
     read the most recent 
     %RY: Maybe freshest is better?
     version. 
    \item Since some readers in progress hold reference to ``stale'' data, the system defers memory reclamation of old data until after each reader in progress has finished executing its read-side critical section. 
\end{enumerate*}

\textit{Remark 2: } RCU read operations can be performed concurrently without any locks, allowing for high concurrency and low contention. This can be particularly useful in systems with multiple subscribers and few publishers. In combination, RCU and Pub-Sub can enable efficient and scalable communication 
% and coordination in concurrent or distributed systems, 
where multiple components need to access and update shared data in a concurrent and asynchronous manner.  

%RY: seems non-central to this paper:
%Therefore, at any point of time, RCU can maintain multiple time-stamped versions of a data item that are concurrently read by readers in the system. 

\subsection{Computing decision updates}\label{DPoverview}
% The decision process as depicted in Fig.~\ref{fig:tc-sys-model} reads the pair of source $1$ and source $2$ updates and generates a decision update based on this pair. 
% Our model assumes that the DP reader can fetch updates from the memory in zero time
% % , we still assume that the DP reader fetches
% albeit
% at some finite average rate such that the combined read request process of all subscribers does not overload the shared memory system.
% We employ a model that captures  the asynchronous operation of the writer and readers of updates in shared memory. 
% 
We view the decision process (DP) reader as one of many subscribers to the updates in the memory system. 
% Because updates are disseminated through the memory system, the writer and the DP reader must operate asynchronously. 
The DP reader becomes aware of fresher updates in the memory only when it chooses to query the memory for a fresh sample of the source update pair. 
We assume a reader can fetch the updates of both source $1$ and source $2$ from memory in negligible time\footnote{This assumption is consistent with RCU reads being lightweight and fast, so that the heavier load is indeed induced by actual decision computation. Further, our model assumes that the DP reader 
% can fetch 
fetches
updates from the memory 
% in zero time
% % , we still assume that the DP reader fetches
% albeit
at some finite average rate such that the combined read request process of all subscribers does not overload the shared memory system.}. 
%From the perspective of the decision process, the shared memory represents a source of updates and 
With this assumption, the DP reader  is an observer that is {\em sampling} the pair of source updates from the memory as a point process. 
Based on this sample, the decision process derives a decision update  which is sent to the monitor, as shown  in Fig.~\ref{fig:tc-sys-model}. The reader process fetches the next sample of update pair from the memory only after the computation in progress is completed. 

When the DP reader's inter-sample times form a renewal process, this is an example of the model of renewal process sampling of updates introduced in \cite{Yates-IT2020}.  In this model, the DP reader generates  an age process $\dpage(t)$  at the input to the DP that is a sampled version of the max-age process $\maxage(t)$ in the memory. Specifically, in the absence of a read, $\dpage(t)$ continues to grow at unit rate. However, if the DP reader makes a read at time $\tau$, then $\dpage(t)$ is reset to $\dpage(\tau)=\maxage(\tau)$. This update pair is then processed by the DP  for a time $T$ so that at time $\tau+T$ a decision update with age $\dpage(\tau)+T$ is delivered to the monitor. The age at the monitor, $\monage(t)$,  is  then reduced to $\monage(\tau+T)=\maxage(\tau)+T$. At this time, the DP reader may choose to fetch a new sample pair from the memory, or it may choose to wait for a time $W$ before fetching the next sample pair. When the DP reader employs non-zero waiting times, we say the DP is using a {\em lazy sampling} policy \cite{Yates-isit2015}. 
Fig.~\ref{fig:tc-age-ev} illustrates the evolution of age processes $\maxage(t), \dpage(t)$, and $\monage(t)$.

\begin{figure}[t]
\centering
\begin{tikzpicture}[scale=0.18] % \linewidth/44cm
\draw [<-|] (0,18) node [above] {age} -- (0,0) -- (16,0);
\draw [|->] (17,0) -- (38,0) node [right] {$t$};
\draw [very thick, dotted] (0,2) -- (4,6) -- (4,4) 
-- (9,9)  -- (9,1) -- (13,5) -- (13,3) -- (15,5);
\draw [very thick, dotted] (18,5) -- (22,9) -- (22,2) -- (27,7) -- (27,4) -- (31,8) -- (31,3) -- (35,7) -- (35,4) -- (37,6) node [right] {$\maxage(t)$};
\draw [very thick, dashed] (0,3) -- (6,9) -- (6,6) -- (14,14) -- (14,4);
\draw [very thick, dashed] (19,9) -- (19,6) -- (23,10) -- (23,3) -- (32,12) -- (32,4) -- (37,9) node [right] {$\hat{y}(t)$};
\draw [very thick] (0,6) -- (3,9) -- (3,6) -- (10,13) -- (10,10) -- (14,14);
\draw [very thick] (18,8) -- (20,10) -- (20,7) -- (29,16) -- (29,9) -- (37,17) node [right] {$\monage(t)$};
\draw 
(0,1) node [left] {$\hat{x}(0)$}
(0,3) node [left] {$\hat{y}(0)$}
(0,5) node [left] {$\hat{z}(0)$}
(6,0) node {$\scriptstyle\blacktriangledown$}
(6,0) node [below] {$\tau_1$} 
(14,0) node {$\scriptstyle\blacktriangledown$}
(14,0) node [below] {$\tau_2$}
(19,0) node {$\scriptstyle\blacktriangledown$}
(19,0) node [below] {$\tau_{i-2}$}
(23,0) node {$\scriptstyle\blacktriangledown$}
(23,0) node [below] {$\tau_{i-1}$}
(32,0) node {$\scriptstyle\blacktriangledown$}
(32,0) node [below] {$\tau_{i}$};
% (3,0.4) -- ++ (0,-0.4) node [below] {$t_1^\prime$}
% (10,0.4) -- ++ (0,-0.4) node [below] {$t_2^\prime$}
% (29,0.4) -- ++ (0,-0.4) node [below] {$t_i^\prime$};
\draw  [|<->|] (0,-2.75) to node [below] {$Y_1$} (6,-2.75);
\draw  [|<->|] (6,-2.75) to node [below] {$Y_2$} (14,-2.75);
\draw  [|<->|] (6,-6) to node [below] {$T_2$} (10,-6);
\draw  [|<->|] (10,-6) to node [below] {$W_2$} (14,-6);
\draw  [|<->|] (19,-2.75) to node [below] {$Y_{i-1}$} (23,-2.75);
\draw  [|<->|] (23,-2.75) to node [below] {$Y_i$} (32,-2.75);
\draw  [|<->|] (23,-6) to node [below] {$T_i$} (29,-6);
\draw  [|<->|] (29,-6) to node [below] {$W_i$} (32,-6);
\draw [|<->|,blue] (25.8, 3) to node [left] {\small $H_i\!$} (25.8, 10);
\end{tikzpicture} \\
\caption{\small Example AoI evolution of the max-age process $\maxage(t)$ at the memory, the sampled max-age process $\hat{y}(t)$ with {\em lazy sampling} at the input to the DP, and the age process $\monage(t)$ at the monitor. The DP reader samples updates from the memory
at times $\tau_1, \tau_2, \ldots$, marked by $\scriptstyle\blacktriangledown$. $Y_i$ is the sampling period for sample $i$, $T_i$ is the computation time for decision update based on sample $i-1$, and $W_i$ is the waiting time to get the $i^{th}$ sample.}
\label{fig:tc-age-ev}
\end{figure}

% \subsection{Paper Overview}
%RY lovely but noinessential sentences: 
%The aforementioned system characteristics dictate that the data held by decision process might be stale compared to data in the memory. 
% In order to understand the complex interactions between the age of source updates in the memory and its effect on the timeliness of decision updates, 


 %In fact, we show that the contribution of the decision process to the average AoI at the monitor is additive to the AoI of the max-age process $\maxage(t)$. 
%In some ways, this is not surprising, since the sampling and computation by the DP has no influence on the  age of updates in the memory. Moreover, 


% Particularly, \Thmref{avg-dpage} in section \ref{sec:decision-processing} provides the average age of sampled age process $\dpage(t)$ when
% sampling is through a renewal process in which the renewal points mark the time instances when the reader fetches a new pair from the memory.
% This is followed by studying the delayed-delivery process $\monage(t)$ induced from delivering the decision updates to the monitor. In this regard, \Thmref{avg-monage-wait} provides the analytical expression for the average age of decision updates.
% % delivery of computed decision updates to the monitor induces a delayed-delivery process $\monage(t)$ at the monitor.
% We note a crucial insight from the analysis: at the input of decision process, because the reader process samples randomly and independent of the way updates arrive in the memory i.e. independent of the evolution of $\maxage(t)$ age process, we arrive at the expression for sampled age process using the backward excess of the renewal process. 
% At the output, we are generating decision updates,
% % the samples occur
% with the evolution of the input. 

% \MOVE{The sources are assumed to be equally important in that the minimum age of each source update in memory is desired.} 

\subsection{Paper Overview and Contributions} 
We 
%systematically 
divide our AoI analysis into two stages: 1) We analyze the average age of updates in shared memory.  2) We analyze the additional delay induced by the  decision process computations.  
First, section \ref{sec:SHS} presents a stochastic hybrid system (SHS) evaluation of  the update age processes in the memory. For the system with sources $i=1,2$, we derive the stationary expected ages $\E{x_i(t)}$ as well as the expected age of the {\em max-age} process $\maxage(t)=\max(x_1(t),x_2(t))$, . 

In section~\ref{sec:decision-processing}, stage two of our analysis, we 
evaluate the age $\monage(t)$ of the decision update process at the monitor. 
The decision process is said to be {\em sampling} the source updates from the memory as it holds a sample of updates that were written to the memory.
Even though the sampling and computation of the DP makes no attempt to use the age of its sampled updates to optimize its operation, 
we show that a lazy sampling policy  will be able to reduce  $\monage(t)$. Here we will see that analysis of $\monage(t)$ is separable from the prior SHS analysis of the max-age process $\maxage(t)$  in the shared memory. In particular, the AoI reduction afforded by lazy sampling can be applied to any stationary update age process that is sampled by the DP. 

\section{Age of source updates in the memory} \label{sec:SHS}
Let $U_{i,1}, U_{i, 2} \ldots$ be the sequence of source $i$ update publication times. At any time $t$, $N_i(t)$ source $i$ updates have been published in the memory, and the most recent update is published at time $U_{i, N_i(t)}$.  It follows that the source $i$ update process has age
%\begin{equation}
 $   x_i(t) = t - U_{i, N_i(t)}$
%\end{equation}
in the memory.
Under this model, the update age $x_i(t)$ is reset to the write time $W \sim \exp(\mu)$ when it is published at time $U_{i,N_i(t)}$. 
% \MOVE[Maybe discard?]{The service facility provided by the writer has no waiting room and an update if accepted by the system must only be in service. In this regard, the writer may be idle, or it may be servicing a source $1$ or a 
% source $2$ update.}
% \RY{Text deleted ... better to just go straight to the math.}
% Two cases arise after the writer finishes writing and becomes idle. 
% In the first case, when the writer finishes writing source $1$ update, the age of source $1$ update is reset at the memory. By virtue of this, the age of source $1$ becomes less than age of source $2$ update in the memory. The second case is similar to first with the age of source $2$ update being less than source $1$ after the writer publishes source $2$ update.
%RY max-age already defined
%Let $\maxage(t) = \max(x_1(t), x_2(t))$, the {\em max-age process}, denote the  age of the oldest update in the memory. 
% We refer to $\maxage(t)$ as the {\em max-age process}. 
When the writer writes a fresh source $i$ update at time $t'$, the max-age process $\maxage(t)$ is reset to $\hat{x}(t') = x_j(t')$, with $j \neq i$. In the following, we use a Stochastic Hybrid System (SHS) to capture the evolution of update age processes in the memory.


\subsection{SHS Overview}
To evaluate AoI of source updates, we use a Stochastic Hybrid Systems (SHS)~\cite{hespanha2006modelling} approach,  a technique introduced for AoI evaluation in \cite{yates2018ToIT} and since employed in AoI evaluation of a variety of status updating systems \cite{Yates-aoi2018,Farazi-KB-aoi2018,Maatouk-AE-aoi2019,Kaul-Yates-isit2018priority,Maatouk-AE-ToN2020,Yates-IT2020,Moltafet-LC-CommLetters2021,Moltafet-LC-ISWCS2021}.
A stochastic hybrid system has a state-space with two components -- a  discrete component $q(t) \in \Qcal = \{0, 2, \ldots, M\}$ that is a continuous-time finite-state Markov Chain and a continuous component $\xv(t) = [x_0(t), \ldots , x_n(t)] \in \R^{n+1}$. In AoI analyses using SHS, each $x_{j}(t) \in \xv(t)$ describes an age process of interest. Each transition $l \in \Lcal$ is a directed edge $(q_l, q'_l)$ with a transition rate $\laml$ in the Markov chain. The age process vector evolves at a unit rate in each discrete state $q \in \Qcal$, i.e., $\frac{d\xv}{dt} = \dot{\xv}(t) = \onev$. A transition $l$ causes a system to jump from discrete state $q_l$ to $q_l^\prime$ and resets the continuous state from $\xv$ to $\xv'$ using a linear transition reset map $\Amat_l \in \{0,1\}^{(n \times n)}$ such that $\xv' = \xv \Amat_l$. For simple queues, examples of transition reset mappings $\set{\Amat_l}$  can be found in \cite{yates2018ToIT}.

For a discrete state $\qbar \in \Qcal$, let %$\Lcal_{\qbar}$ and $\Lcal'_{\qbar}$ be sets of incoming and outgoing transitions, i.e.
%\begin{subequations}
\begin{align}
    \Lcal_{\qbar} &= \{l \in \Lcal : q'_l = \qbar \}, &
    \Lcal'_{\qbar} &= \{l \in \Lcal : q_l = \qbar \}.
\end{align}
denote the respective sets of incoming and outgoing transitions.
%\end{subequations}
%Specifically, average age of an age process of interest is $\age_i = \E{x_i} = \sum_{\qbar\in\Qcal} \vbar_{\qbar i}$ and 
Age analysis using SHS is based on the  expected value processes $\set{\vv_q(t)\colon q\in \Qcal}$ such that 
%\begin{align}
 $\vv_q(t)%=\rvec{v_{q1}(t)&\cdots & v_{qn}(t)}
 %&
 =\E{\xv(t)\delta_{q,q(t)}}$,
 %\nn&
 %=\rvec{\E{x_0(t)\delta_{q,q(t)}} & \cdots&\E{x_n(t)\delta_{q,q(t)}}},
%\end{align}
 with $\delta_{i,j}$ denoting the Kronecker delta function. For the SHS models of age processes considered here, each $\vv_q(t)$ will converge to a fixed point $\vvbar_q$. The fixed points $\set{\vvbar_q\colon q\in\Qcal}$ are the solution to a set of age balance equations. The following theorem provides a simple way to calculate the age-balance fixed point and then the average age. 
\begin{theorem}\thmlabel{AOI-SHS}
\cite[Theorem~4]{yates2018ToIT}
If the discrete-state Markov chain $q(t)\in\Qcal=\set{0,\ldots,M}$ is ergodic with stationary distribution 
$\bar{\piv}=\rvec{\bar{\pi}_0 &\cdots &\bar{\pi}_M}>0$ and there exists a 
non-negative vector $\vvbar=\rvec{\vvbar_0&\cdots&\vvbar_M}$
such that 
\begin{align}
\bar{\vv}_{\qbar}\sum_{l\in\Lcal_{\qbar}}\laml &=\onev[]\bar{\pi}_{\qbar}+ \sum_{l\in\Lcal'_{\qbar}}\laml \bar{\vv}_{\ql}\Amat_l,\quad \qbar\in\Qcal,\eqnlabel{AOI-SHS-v}
\end{align}
then
the average age vector is 
$\E{\xv}\!=\!{\displaystyle \limty{t}}\E{\xv(t)}\!=\!
\sum_{\qbar\in\Qcal} \vvbar_{\qbar}$.
\end{theorem}


\subsection{SHS Analysis of Age in Shared Memory}
\begin{figure}[t]
$\setlength{\extrarowheight}{0.5mm}
\begin{array}[b]{cccc}
l & q_l\to q^\prime_l & \laml &  \xv\Amat_l \\ \hline
1 & 0_1\to1 & \lam_1 & \rvec{0, x_1, x_2, \hat{x}} \\
2 & 1\to0_2 & \mu & \rvec{x_0, x_0, x_2, x_2} \\
3 & 0_2\to1 & \lam_1 & \rvec{0, x_1, x_2, \hat{x}} \\
4 & 0_2\to2 & \lam_2 & \rvec{0, x_1, x_2, \hat{x}} \\
5 & 2\to0_1 & \mu & \rvec{x_0, x_1, x_0, x_1} \\
6 & 0_1\to2 & \lam_2 & \rvec{0, x_1, x_2, \hat{x}}
\end{array}$
\hfill
\begin{tikzpicture}[->, >=stealth', auto, semithick, node distance=2cm]
\tikzstyle{every state}=[fill=none,draw=black,thick,text=black,scale=0.85]
\node[state]    (0)                     {$0_1$};
\node[state] (1)[above right of=0] {$1$};
\node[state] (2)[below right of=0] {$2$};
\node[state] (3)[below right of=1] {$0_2$};
\path
% (0) edge[loop left] node{\small $\gam$} (0)
(0) edge[bend left=25, above] node{\small $\lam_1$}  (1)
(0) edge[bend right=25, below] node{\small $\lam_2$} (2)
% (1) edge[loop above] node{\small $\gam$} (1)
(1) edge[bend right=25, below] node{\small $\mu$} (3)
% (3) edge[loop right] node{\small $\gam$} (3)
(3) edge[bend right=25, above] node{\small $\lam_1$} (1)
% (2) edge[loop below] node{\small $\gam$} (2)
(3) edge[bend left=25, below] node{\small $\lam_2$} (2)
(2) edge[bend right = 25,above] node{\small $\mu$} (0);
\end{tikzpicture}
%\end{tabular}
\caption{The SHS transition/reset maps and Markov chain for the update age in the shared memory.}
\label{fig:simplified-tc-mc}
\end{figure}
% \V{Move the SHS Analysis from Appendix to this section}
% We now describe SHS framework for our system model shown in Fig.~\ref{fig:tc-sys-model}.
The age of updates in a shared memory system 
% with RCU read-write access and 
with bufferless service at the writer  can be described by the SHS Markov chain and table of state transitions shown in Fig.~\ref{fig:simplified-tc-mc}.  The continuous age state vector 
%for our SHS framework is $\xv \in \R^4$, with 
is $\xv = \rvec{x_0, x_1, x_2, \hat{x}}$, where 
$x_0$ is the age of the update being written; $x_i$, $i=1,2$,  is the age of the source $i$ update in memory; and
$\hat{x}=\max(x_1,x_2)$. 
% Fig.~\ref{fig:simplified-tc-mc} shows a 
% SHS 
%The Markov Chain 
% for the 
The discrete state is $\Qcal = \set{0_1, 0_2, 1, 2}$. At time $t$, the system  is in state $0_i$ if the writer is idle and the oldest update belongs to source $i$. State $i\in \set{1, 2}$ corresponds to the writer writing source $i$ update.

We now describe SHS transitions enumerated in the table in Fig.~\ref{fig:simplified-tc-mc}. 
For each collection of transitions, we focus on the age state components 
% $x_i(t), i \in \set{0, 4}$ 
that change.
\begin{itemize}
\item {$l=1,3,4,6$:} In system idle states $0_1$ and $0_2$, the writer receives a new source update and initiates a new write mechanism. $x_0^\prime = 0$ as the writer receives a fresh update, and $x_1^\prime , x_2^\prime, \hat{x}^\prime$ are unchanged as the update is not yet written to the memory. 
% and the decision process is still busy computing a decision update.
% Notice that this transition is source agnostic.

\item {$l=2, 5$:} The writer finishes writing and publishes a new source update. 
    \begin{description}
        \item[$l=2$]: the writer publishes source $1$ update: $x_1^\prime = x_0$ as the age of source $1$ update in the memory is reset to just written update. The source $2$ update becomes the oldest update in the memory; hence, 
        $\hat{x}^\prime = x_2$. 
        \item[$l=5$]: The writer publishes source $2$ update: $x_2^\prime = x_0$, the source $1$ update becomes the oldest update, and $\hat{x}^\prime = x_1$.
    \end{description} 
\end{itemize}
\noindent For the SHS analysis, we  employ  the normalized rates \begin{equation}
    \rho_1={\lam_1}/\mu,\quad \rho_2={\lam_2}/{\mu}.
%\sig_{\RCU}=\mu_{\RCU}/\mh,\quad \sig_{\RWL} = \mu_{\RWL}/\mh.
%RY defer intro of sig versions for RCU and RWL
\end{equation} 
We note that $\rho = \rho_1 + \rho_2$ is the total offered load of source updates being written to the memory. 
% Similarly, $\sig$ is the normalized decision update generation rate.
The Markov chain in Fig.~\ref{fig:simplified-tc-mc} has stationary probabilities $\piv$
%$\piv =\rvec{\pi_0 & \pi_1 & \pi_2 & \pi_3 & \pi_4}$ 
with normalization constant $C_\pi$ given by
\begin{subequations}
\begin{align}
    \piv &=  \rvec{\pi_{0_1} & \pi_1 & \pi_2 & \pi_{0_2}}
    = C_\pi^{-1} \rvec{\rho_2 / \rho & \rho_1 & \rho_2 & \rho_1/ \rho},\\
%C_\pi &= \bt\rh(1+\bt+\rh) + \sig(1+\bt)(1+\rh)(\bt+\rh)\nn  &\qquad+\sig^2(1+\bt)(1+\rh).
C_\pi &= 1+\rho.
\end{align}
\end{subequations}
% We now apply \Thmref{AOI-SHS} to  the SHS system in Fig.~\ref{fig:simplified-tc-mc}. 
With the shorthand notation 
\begin{equation}
    \lam=\lam_1+\lam_2,
    \end{equation}
% and$\piv$, 
% \RY{What's the difference between $\pi$ and $\pibar$? Also why use $\lam$ instead of just $\lam$?}
we now use \Thmref{AOI-SHS} to solve for
\begin{align}
    \vvbar &= \rvec{\vvbar_{0_1}&\vvbar_{1}&\vvbar_{2}&\vvbar_{0_2}},
\end{align}
where 
% the decision update age process $\hat{\xv}(t)$ has age balance fixed points
$\vv_q=\rvec{v_{q0}&v_{q1}&v_{q2}&v_{q3}}, \forall q \in \Qcal$. 
%RY Q was already enumerated $= \set{0_1, 1, 2, 0_2}$.
This yields
\begin{subequations}
\eqnlabel{simplified-tc-shs-eqn}
\begin{align}
\lam \vvbar_{0_1} &= \onev[]\pibar_{0_1} + \mu \vvbar_2 \Amat_5, \\
\mu\vvbar_1 &= \onev[]\pibar_1 + \lam_1 \vvbar_{0_1} \Amat_1 + \lam_1 \vvbar_{0_2} \Amat_3, \\
\mu\vvbar_2 &= \onev[]\pibar_2 + \lam_2 \vvbar_{0_1} \Amat_6 + \lam_2 \vvbar_{0_2} \Amat_4, \\
\lam \vvbar_{0_2} &= \onev[]\pibar_{0_2} + \mu \vvbar_1 \Amat_2.
\end{align}
\end{subequations}
%%%%%%
%% See explanation in comments after Conclusion
% \enlargethispage{-1.4cm} 
%%%%%%
We can now use \Thmref{AOI-SHS} to calculate the AoI of source $i$ update in the memory as $\E{x_i} = v_{0_1, i} + v_{0_2, i} + v_{1, i} + v_{2, i}$ for $i \in \{1, 2 \}$ as well as $\E{\hat{x}}= v_{0_1, 3} + v_{0_2, 3} + v_{1, 3} + v_{2, 3}$. 
% With 
% additional
Some
algebra
% , \Thmref{age-exp} can be verified.
% An SHS analysis using \Thmref{AOI-SHS} 
yields the following theorem.
% Two cases arise after the writer finishes writing and becomes idle. 
% % Consider when a type $1$ update is in service. 
% In the first case, when the writer finishes writing source $1$ update, the age of source $1$ update is reset at the memory. By virtue of this, the age of source $1$ becomes less than age of source $2$ update in the memory. The second case is similar to first with the age of source $2$ update being less than source $1$ after the writer publishes source $2$ update.
% % The discrete states of our SHS framework accounts for these two cases. 
% Fig.~\ref{fig:tc-mc} shows a Markov Chain for the discrete state $\Qcal = \set{0_1, 0_2, 1, 2}$.
% % Therefore, the discrete states of the Markov Chain will be $\Qcal = \set{0_1, 0_2, 1, 2}$, 
% At time $t$, the system  is in state $0_i$ if the writer is idle and $\Delta_i(t)=\max(\Delta_1(t),\Delta_2(t))$.
%provides analytical expressions for the average ages  $\E{x_i(t)}$ 
 %of source $i$ updates and of the max-age process 
 %and $\E{\maxage(t)}$ in shared memory. 
\begin{theorem}\label{thm:age-exp}
% For load parameters $\rho_i = \lam_i/\mu$ such that $\rho = \rho_1 + \rho_2$,
\mbox{}
\begin{enumerate}
\item[(a)] Source $i$ updates in the memory have average age
\begin{align}
    \E{x_i} = \frac{1}{\mu}\left(\frac{1+\rho}{\rho_i} + \frac{\rho}{1+
    \rho}\right).
    % \frac{1}{\mu \rho_i} (1+\rho_j + (1+ C) \rho_i),
    \eqnlabel{source-i-age}
\end{align}
% where, $C = (\rho_i + \rho_j) / (1+\rho_i+\rho_j)$.
% \vspace{4mm}
\item[(b)] The max-age process $\hat{x}(t)=\max(x_1(t), x_2(t))$ in the memory has average age
% \begin{align}
%     \E{\hat{x}} = \frac{\rho_1^2(1+\rho_1) + \rho_2(B(\rho_1 + \rho_2) + \rho_2^2)}{\mu \rho_1 \rho_2 (\rho_1 + \rho_2)},
%     \eqnlabel{max-age}
% \end{align}
% where, $B = 1 + (2+C)\rho_1$. 
% \RY{I did some algebra and got the symmetric expression
\begin{align}\eqnlabel{Emaxage}
\E{\xhat}&=\frac{(1+\rho)^2(\rho_1^2+\rho_1\rho_2+\rho_2^2)+\rho^2\rho_1\rho_2}{\mu\rho(1+\rho)\rho_1\rho_2}.
\end{align}
% }
\end{enumerate}
\end{theorem}
% \noindent The proof, including explanations of the SHS state transitions, appears in the Appendix. While the bulk of this proof is the enumeration of the age balance equations \eqnref{AOI-SHS-v}, some readers may be interested in how the max-age $\maxage(t)=\max(x_1(t),x_2(t))$ is incorporated as one of the continuous variables in the SHS framework. 

\begin{figure}[t]
\centering
\includegraphics{MATLAB/maxageplotnew.pdf}
\caption{Average age of max-age process $\maxage(t)$ in the memory. For a fixed updating load, we vary  $\alpha$ with $\rho_1=\alpha\rho$  and $\rho_2=(1-\alpha)\rho$.
}
\label{fig:maxageplot}
\end{figure}
\noindent
Not surprisingly, the expected max-age $\E{\maxage}$ is symmetric in the  load parameters $\rho_1$ and $\rho_2$. However, since  the formula \eqnref{Emaxage}  is somewhat opaque, a plot of $\E{\maxage}$ appears in Fig.~\ref{fig:maxageplot}. A possibly non-obvious observation from the figure is that increasing the overall updating load $\rho$ generally improves the average max-age because the writer queues no updates. The figure also reveals that the average max-age is penalized by asymmetry in the update rates of the individual sources. This is in part because a source that updates slowly will have high age and thus cause the max-age to be large.  However, it is also true that with asymmetric loads, the high rate  source  will cause updates of the low rate source to be discarded at the writer. Because the writer is non-selective in offering service, it may be performing updates for the high rate source even when the age of that source is already low.  

 






\section{Age of decision updates} \label{sec:decision-processing}
In section~\ref{DPoverview}, we observed that the DP reader is {\em sampling} the source updates from the memory as a point process. 
%
% This phenomenon requires us to define the age process induced by sampling at the decision process:  
% %
% \begin{definition} \label{def:sampling}
% An age process $y(t)$ {\em samples} the age update process $x(t)$ of a source if
% \begin{enumerate}
% \item There is a point process that marks the time instances that the monitor fetches the current update of the source.
% \item When source update is fetched at time $t^\prime$, $y(t')$, the age of the source at the monitor is reset to $y(t') = x(t')$. 
% \item In the absence of a new read, $y(t)$ grows at unit rate. 
% \end{enumerate}
% \end{definition}
%
% \RY{I tried to re-write Definition 1 in a generic way so as to not give the impression that we will be sampling $x_1(t)$ and perhaps $x_2(t)$ independently. In fact, I think we are sampling the max-age process $\hat{x}(t)$ in which we are really collecting samples of both $x_1$ and $X_2$. Is that right?} 
% In our system, the shared memory is the source and the reader, on behalf of the decision process,  is the  monitor that samples the stationary max-age process $\xhat(t)$. In particular, when the reader fetches an update, it fetches the current updates of both sources $1$ and $2$. At the input to the decision process, we denote sampled max-age process by $\hat{y}(t)$. 
%
In particular, we assume the inter-sample times $Y_1,Y_2,\ldots$ that are
i.i.d continuous random variables identical to $Y$. In this case, the update sample times form a renewal process, and in the parlance of \cite{Yates-IT2020}, the update age process $\dpage(t)$ is sampling the max-age update process $\maxage(t)$ in the shared memory.  

\subsection{Average Age at the Decision Process}
At %an arbitrary 
time $t$, the most recent read from memory occurred at time $t-Z(t)$. That is,  $Z(t)$ is the age of the sampling renewal process.  When the renewal process is in equilibrium,  $Z(t)$ is stationary with first moment 
\cite[Theorem 5.7.4]{Gallager2013stochastic}
% \begin{equation}
%     f_Z(z) = \frac{P[Y > z]}{E[Y]}, \quad z \geq 0, \eqnlabel{bcwexc-pdf}
% \end{equation}
\begin{equation}
\E{Z} = \frac{\E{Y^2}}{2\E{Y}}. \eqnlabel{bcwexc-mom}
\end{equation}
Next, following the approach in \cite{Yates-IT2020}, we observe that the DP reader does not fetch any update in the interval $(t-Z(t),t]$. Hence, at time $t$, the update age $\dpage(t)$ satisfies
\begin{equation}
\dpage(t) = \maxage(t - Z(t)) + Z(t). \eqnlabel{yhat-t}
\end{equation}
Further, $Z(t)$ is independent of $\maxage(t)$ because the inter-sample times $Y_i$ are independent of the age processes in the shared memory. Thus 
 stationarity of $\E{\maxage(t)}$ implies 
 $\E{\maxage(t-Z(t)}=\E{\maxage(t)}=\E{\maxage}$. It then follows from \eqnref{bcwexc-mom} and \eqnref{yhat-t} that $\dpage(t)$ has expected value\footnote{A stronger distributional result is derived in \cite[Theorem~6]{Yates-IT2020} that is not needed for the average AoI analysis here.}  
\begin{equation}
\E{\dpage} = \E{\maxage} + \E{Z}=\E{\maxage} + \frac{\E{Y^2}}{2\E{Y}}. \eqnlabel{Eyhat-Y}
\end{equation}

% \V{however, I have not yet explicitly stated how and why $\maxage(t)$ is stationary.} \RY{The SHS analysis does confirm that the the age processes are stationary in the mean. That is, if the system runs for a long time (or was started a long time ago) then $E{x(t)}$ is constant. Similarly the MGF also converges. This is good enough for the average AoI analysis of the computation.} 

%RY This paragrpah might be better than what's in the Overview
% \V{Changed the name lazy computation to lazy ssampling}
\subsection{Average Age at the Monitor: Lazy Sampling} 
When the DP reader samples the shared memory, the DP then computes a decision update based on this sample. On delivery of a decision update to the monitor, the update age $\monage(t)$ is reset to the age of the oldest source update that was read and used to compute the decision update. This means that an arrival of decision update at the monitor at time $t$ resets $\monage(t)$ to $\dpage(t)$. 

In this work, we assume that the decision computation times are i.i.d continuous random variables $T_1, T_2, \ldots$, each identically distributed to $T$. We will consider a DP that performs {\em lazy sampling}: after delivering the computation to the output monitor, the DP reader waits for a random time $W$ before reading again. The alternative
    to being lazy is the {\em zero-wait} policy, a special case of lazy when $W=0$.
%While zero-wait is a seemingly sensible policy in that we don't waste available processor cycles.

%\subsection{Lazy computation}
% We first consider a general scenario where the decision process after delivering the computation to the monitor waits for some random time denoted by $W$.
Fig.~\ref{fig:tc-age-ev} depicts the evolution of the max-age process $\maxage(t)=\max(\oneage(t), \twoage(t))$, the status-sampling process $\dpage(t)$, and the age at the monitor $\monage(t)$.  with i.i.d inter-sample intervals $Y_1, Y_2, \ldots$ such that samples are taken at times $\tau_i = \sum_{j=1}^i Y_j$. 

Under lazy sampling, we admit the possibility that the $i$th computation time $T_i$ and the $i$th waiting time $W_i$ are correlated. However, in order for the $\dpage(t)$ process to be sampling the shared memory with independent inter-sample times $Y_i=T_i+W_i$, we require that the pairs $(T_1,W_1),(T_2,W_2),\ldots$ to be i.i.d., identical to $(T,W)$. Under this assumption, it follows directly from \eqnref{Eyhat-Y} that the average update age at the input to the DP is
\begin{equation}\eqnlabel{DPageTW}
\E{\dpage} = \E{\maxage}+\frac{\E{(T+W)^2}}{2\E{T+W}}.
\end{equation}

Curiously, \eqnref{DPageTW} reveals that the problem of minimizing the average age at the input to DP appears to be isomorphic to the timely updating problem that was originally formulated in \cite{Sun-UBYKS-IT2017UpdateorWait,Yates-isit2015}, where the suboptimality of zero-wait policies was first identified. However, in this system, our objective is not to minimize $\E{\dpage}$ but rather to minimize the average age $\E{\monage}$ at the monitor.  Since $\monage(t)$ is penalized by the waiting time $W$, choosing $W$ to minimize $\E{\dpage}$ may not be good for $\E{\monage}$.
%the end-to-end AoI. 
Fortunately, the following claim verifies this is not the case. 

% The computation from sample $i$ is delivered to the monitor at $t_i^\prime = \tau_{i-1} + T_i$, hence $\monage(t_i^\prime) = \dpage(t_i^\prime)$. A new sample $i + 1$ is fetched at time $\tau_i = \tau_{i-1} + T_i + W_i$ by decision process, therefore, $\dpage(\tau_i) = \maxage(\tau_i)$.

% We provide a general analysis for the age of decision update where computation time and waiting time can have arbitrary distributions. \Thmref{avg-monage-nowait} is then shown as a special case of \Thmref{avg-monage-wait}.

% \begin{theorem} \thmlabel{avg-monage-wait}
%     If $\maxage(t)$ is the ergodic max-age process at the memory with mean $\E{\maxage(t)}$, 
%     then for i.i.d 
%     sampling intervals $Y = T+W$, where $T$ and $W$ are 
%     %independent
%     random decision computation and waiting times, 
%     the age at the monitor $\E{\monage}$ is given as:
%     \begin{align}
%         \E{\monage} &= \E{\maxage(t)} + \frac{\E{Y^2}}{2\E{Y}} + \E{T} \eqnlabel{tmp1} \\
%         &= \E{\maxage} + \frac{1}{2}\left(\E{T} + \E{W} + \frac{\Var{T+W}}{\E{T} + \E{W}}\right) + \E{T}. \eqnlabel{avg-monage-wait}
%     \end{align}
% \end{theorem}
\begin{theorem} \thmlabel{avg-monage-wait}
    If $\maxage(t)$ is a stationary process, 
    then for any  waiting policy such that $W_i$ depends only on $T_i$, the average age 
    %$\E{\monage(t)}$ 
    at the monitor satisfies 
    \begin{align}
        \E{\monage} &= \E{\dpage} +\E{T}.  \eqnlabel{avg-monage-wait}
    \end{align}
\end{theorem}
The proof appears in the Appendix. We observe that \Thmref{avg-monage-wait}  can give one the mistaken impression that $\E{\monage}$ is insensitive to the waiting time $W$. In fact, the theorem says that the waiting time $W$ affects $\E{\dpage}$ and $\E{\monage}$ identically. A hand-waving intuition is that $\monage(t)$ lags $\dpage(t)$ only during the computation time $T$ but, once the computation is complete, $\monage(t)=\dpage(t)$ during any waiting period. 

% \subsection{Zero-wait computation}
% A special case of the lazy policy is that following the decision processing time, the generated decision updates is sent instantly to the monitor and a new pair of source updates is fetched for the next round. 
% Fig.~\ref{fig:tc-age-ev}(a) depicts the evolution of $\maxage(t), \dpage(t)$, and $\monage(t)$ for zero-wait policy.
% At time $\tau_i = Y_1 + Y_2, \ldots + Y_i$,
% a decision update based on sample $i$ is delivered to the monitor, hence $\monage(\tau_i) = \dpage(\tau_i)$. At the same time, the next sample $i+1$ of $\maxage(t)$ is fetched by the decision process, and therefore $\dpage(\tau_i) = \maxage(\tau_i)$,
% as shown in Fig.~\ref{fig:tc-age-ev}(a). With $W=0$ and $Y = T$, it is straightforward to obtain the stationary expected age $\E{\monage(t)}$ using \Thmref{avg-monage-wait}.


% \begin{corollary}

% If $\maxage(t)$ is the ergodic max-age process at the memory with mean $\E{\maxage(t)}$, then for zero-wait computation policy where the sampling intervals $Y$ are identical to computation time $T$, and waiting time $W = 0$, the age at the monitor $\E{\monage}$ is given as:
% \begin{align}
% \E{\monage(t)} &= \E{\maxage(t)} + \frac{\E{Y^2}}{2\E{Y}} + \E{T} \eqnlabel{avg-monage-std} \\
% &= \E{\maxage(t)} + \frac{3}{2}\E{Y} + \frac{\Var{Y}}{\E{Y}}. \eqnlabel{avg-monage-nowait}
% \end{align}
% \end{corollary}


% The analyses in the following subsections reveal a
% a rather surprising result that for both zero-wait and any lazy policy, the average age of the decision update process at the output monitor is given by a common analytical expression
% \begin{subequations}
% \begin{align}
%     \E{\monage(t)} &=  \E{\dpage(t)} + \E{T} \eqnlabel{monage-wrt-dpage}\\
%     & = \E{\maxage(t)} + \frac{\E{Y^2}}{2\E{Y}} + \E{T}.
%     \eqnlabel{avg-monage}
% \end{align}
% \end{subequations}
% This is surprising especially in case of a non-zero waiting policy where 
% even though the sampling is through a renewal process, the inter-update times at the monitor are not necessarily independent. 
% In fact, for the case of zero-waiting policy, 
% during computation time $T_i$, $\monage(t) \neq \dpage(t-T_{i-1}) + T_{i-1}$. 
% Specifically, we observe that $\dpage(t)$ is a renewal sampling process but $\monage(t)$ isn't.

Combining \eqnref{DPageTW} and \eqnref{avg-monage-wait}, we obtain an end-to-end characterization of the average age in the system:
\begin{align}
        \E{\monage} &= \E{\maxage}+\frac{\E{(T+W)^2}}{2\E{T+W}} +\E{T}.  \eqnlabel{avg-monage-wait-final}
    \end{align}
Since the computation time $T$ is given, \eqnref{avg-monage-wait-final} shows that the choice of a waiting function $W$ as a function of $T$ is the same problem formulated in \cite{Sun-UBYKS-IT2017UpdateorWait,Yates-isit2015}. Hence the solution is the same, namely the $\beta$-minimum waiting policy 
\begin{equation}
    W_i = (\beta - T_i)^+,
\end{equation}
where the parameter $\beta$ is chosen by numerical line search. 
With this policy, $T+W=\max(\beta,T)$ and it follows  from \eqnref{avg-monage-wait-final} that the policy achieves end-to-end  average AoI
%age at the output monitor
\begin{equation}
        \E{\monage} = \E{\maxage} + \frac{\E{\max(\beta^2, T^2)}}{2\E{\max(\beta,T)}} + \E{T}. \eqnlabel{beta-min-age}
    \end{equation}
For completeness, the effectiveness of waiting is demonstrated in  section~\ref{num-eval} by some numerical evaluations of the lazy sampling policy. We will see that lazy sampling becomes important when the variance of the computation time $T$ becomes large.  Before presenting these results, we comment on the connection of this lazy sampling model to the lazy updating model in  \cite{Sun-UBYKS-IT2017UpdateorWait,Yates-isit2015}.

In \cite{Sun-UBYKS-IT2017UpdateorWait,Yates-isit2015}, the random variable
$T$ represented the delivery time of a fresh update (say through a network)  to the monitor. Fresh updates were generated at will and $W$ represented the waiting time prior to generating the next fresh update. A key element of this system was the tight coupling of waiting  and update generation.  In this setting, the intuition behind $\beta$-minimum waiting was that if the prior delivery time was small, the age at the monitor would be small and it would be  a waste of network resources to deliver an update when the age reduction afforded by the update would be small. 

In this work, updates are generated by an exogenous process that is beyond the control of the DP. Moreover, because updates are disseminated through a shared memory publication process, the  age processes of updates in shared memory are essentially uncoupled from the update sampling/processing policy implemented by the DP. In particular, any time the DP reader fetches a sample pair from the memory, the update age of that pair has expected value $\E{\maxage}$, which is just the average age in the shared memory. Nevertheless, even though  DP operation is uncoupled from the age process in shared memory, the $\beta$-minimum waiting policy is effective. In particular, it reduces the expected value of $\dpage(t)$, the age process at the input to the DP. What is happening is that the waiting policy mitigates the deleterious effect of high-variance computation times $T$ on the sampling policy at the DP reader.
We note that \Thmref{avg-monage-wait} went unrecognized in \cite{Sun-UBYKS-IT2017UpdateorWait,Yates-isit2015}. Specifically, \Thmref{avg-monage-wait} shows that no matter what policy is used, the output always lags the input by $\E{T}$ in terms of average age.
% with high-variance inter-sample times. 

% However, various studies in the AoI literature suggest that the zero-wait policy  is sub-optimal \cite{Yates-isit2015, Sun-UBYKS-infocom2016UpdateorWait}. Since the decision process can insert waiting times before fetching the next sample from the memory, how to optimally sample to maximize data freshness at the monitor bears examination. In this regard, we consider and extend the setup introduced in \cite{Yates-isit2015} as follows.

% \subsection{$\beta$-minimum policy}
% In the system with waiting regime, the inter-sample time is $Y_i = T_i + W_i$. By choosing $W_i$ as a function of $T_i$, $Y_i$ are i.i.d random variables. We fix the sampling rate:
% \begin{equation}
%     \gamma = \frac{1}{\E{T} + \E{W}} \eqnlabel{lam}
% \end{equation}
% From \Thmref{avg-monage-wait}, the average age at the output monitor is:
% \begin{align}
%     \E{\monage(t)} &= \E{\maxage(t)} + \frac{\E{Y^2}}{2\E{Y}} + \E{T} \\
%     &= \E{\maxage(t)} + \frac{\E{(T+W)^2}}{2\E{(T+W)}} + \E{T} \\
%     &= \E{\maxage(t)} + \frac{\gamma}{2} \E{(T+W)^2} + \E{T}.
% \end{align}
% As suggested in \cite{Yates-isit2015}, for fixed $\gamma$, 
% % \cite{Yates-isit2015} says that
% we can minimize $\E{\monage(t)}$ by choosing $W$ as a deterministic function of $T$, $W = z(T)$. Therefore, we fix $\gamma$, treat \eqnref{lam} as a constraint on $\E{W}$ and minimize $\E{\monage(t)}$ by minimizing $\E{(T+W)^2}$. The solution of this functional minimization is given by $\beta$-minimum policy in \cite[Theorem 1]{Yates-isit2015}. Specifically, we have a $\beta$-minimum policy of the following Theorem.
% \begin{theorem}
%     For a given sample rate $\gamma \leq 1/\E{T}$, the age at the monitor is minimized by choosing
%     $$W_i = (\beta - T_i)^+$$
%     such that $\beta$ satisfies,
%     \begin{equation}
%         \E{\max(\beta, T)} = \frac{1}{\gamma} \eqnlabel{beta-const}
%     \end{equation}
%     Such a policy achieves decision update age
%     \begin{equation}
%         \E{\monage(t)} = \E{\maxage(t)} + \frac{\gamma}{2}\E{\max(\beta^2, T^2)} + \E{T}. \eqnlabel{beta-min-age}
%     \end{equation}
% \end{theorem}
% \begin{proof}
% We let $\alpha = 1/\gamma - \E{T}$. To ensure non-negativity, we write $W = z(T) = w^2(T)$ and optimize over $w(t)$. This means that $\E{w^2(T) = \alpha}$. With $\beta$ as the Lagrange multiplier:
% \begin{align}
% L(w) &= \E{(T+w^2(T))^2} + 2\beta\E{\alpha - w^2(T)} \\
% &= \int_0^\infty ((t+w^2(t))^2 + 2\beta(\alpha - w^2(t)))f_T(t) dt\
% \end{align}
% We observe the Euler-Lagrange differential equation with
% \begin{equation}
% L(w) = \int_0^\infty F(t, w, w^\prime) dt
% \end{equation}
% where,
% $$F(t,w,w^\prime) = (t+w^2(t))^2 + 2\beta(\alpha - w^2(t)) f_T(t)$$
% $L(w)$ has a stationary value if the Euler-Lagrange differential equation
% \begin{equation}
%     \frac{\partial F}{\partial w} - \frac{d}{dt} \frac{\partial F }{\partial w^\prime} = 0
% \end{equation}
% is satisfied.
% Since, $F(t,w,w^\prime)$ is not a function of $w^\prime$, the Euler-Lagrange differential equation reduces to
% \begin{equation}
% \frac{\partial F}{\partial w} = 4((t + w^2)w - \beta w) f_T(t) = 0
% \end{equation}
% which implies $(w^2(t) - \beta + t)w(t) = 0$, which gives us $z(t) = w^2(t) = (\beta - t)^+$. By choosing $\beta$ such that $W = z(T)$ satisfies \eqnref{lam}, \eqnref{beta-const} and \eqnref{beta-min-age} follow.
% \end{proof}



%RY I just did some editing on this paragraph. I can't tell whehter you or I started writing this.
% In the Lazy is Timely, we were really exploiting the fact that we were choosing the wait relative to age of update process, for instance in lazy is timely, the beta policy is beta - previous age. The paper relied on the fact that at the end of the computation, you start with a fresh update. Here you don't get to choose fresh update, you simply get to choose a delay until you get a random (possible correlated) sample of the update. But effectively, the source updates are recorded out of your control, there is still this same optimum way to do the computation phase to mitigate the AoI penalty associated with computation.

\section{Numerical Evaluation}\label{num-eval}
\begin{figure}[t]
\centering
\begin{tabular}{c}
\includegraphics{MATLAB/betaminexp.pdf} \\
\qquad \textbf{(a)} $T\sim\exp(1)$ \\[3mm]
\includegraphics{MATLAB/betaminlognormal.pdf} \\%[3mm]
\qquad \textbf{(b)} $T\sim$ Log-Normal($1, e^{\sig^2} - 1$)
\end{tabular}

\medskip

\caption{Average age at the monitor vs the sampling rate $\gamma$ for the $\beta$-minimum policy for different distribution of computation time $T$. Total offered load by source updates is $\rho=1$,
with $\rho_1 = \rho_2 = 0.5$. 
Notice that $\gamma = 1$ is the zero-wait computation policy.}
\label{fig:beta-min-plot}
\end{figure}





% \begin{figure}[t]
% \centering
% \includegraphics{MATLAB/betaminexp.eps}
% \caption{Average age at the monitor vs the sampling rate $\gamma$ for the $\beta$-minimum policy and $T \sim \exp(1)$ computation time with $\rho_1 = \rho_2 = 0.5$. Notice that $\gamma = 1$ is a zero-wait sampling policy.}
% \end{figure}

% \begin{figure}[t]
% \centering
% \includegraphics{MATLAB/betaminlognormal.eps}
% \caption{Average age at the monitor vs the sampling rate $\gamma$ for the $\beta$-minimum policy. Computation time $T$ is log-normal with $\E{T} = 1$, and $\Var{T} = e^{\sig^2} - 1$. Notice that $\gamma = 1$ is a zero-wait sampling policy.}
% \end{figure}

% \begin{figure}[t]
% \centering
% \includegraphics{MATLAB/betamin_sig_gam1.eps}
% \caption{$b = -\sig^2/2, \gamma = 1$}
% \end{figure}
In this section, we  examine some numerical examples of the performance $\beta$-minimum waiting policy, simply to remind the reader of the benefits of waiting. 
Fig.~\ref{fig:beta-min-plot} illustrates age performance with respect to variance in the computation time with probability distributions $\exp(1)$ (Fig.~\ref{fig:beta-min-plot}(a)), and Log-Normal($1, e^{\sig^2} - 1$) (Fig.~\ref{fig:beta-min-plot}(b)). 
The log-normal distributed computation times $T$ has PDF \cite{yates2014probability},
\begin{align}
f_T(t) &= %\begin{cases} %RY eveyrone knows T is non-negative
\frac{e^{-(\ln(t) - b)^2/2\sig^2}}{\sqrt{2\pi}\sig t}, \qquad t > 0,
%0 & \text{otherwise}.
%\end{cases}
\end{align}
with free parameters $b$ and $\sig > 0$.
In our numerical evaluations, 
we consider a given distribution on $T$ such that the computation time is normalized to $\E{T}=1$. In this regard, for Log-Normal distribution,
for each $\sig$, we set $b = -\sig^2/2$ so that $\E{T} = 1$. By varying $\sigma$, we vary $\Var{T}=e^{\sigma^2}-1$.
These numerical results are largely similar to those in \cite{Sun-UBYKS-IT2017UpdateorWait,Yates-isit2015}.
In particular, the results remind the reader that zero-wait becomes increasingly sub-optimal when the computation time $T$ has high variance. 
% In each example, we consider a given distribution on $T$ such that the computation time is normalized to $\E{T}=1$.  
The choice of $\beta$ specifies a sampling rate  
\begin{equation}
\gamma = \frac{1}{\E{T} + \E{W}} =\frac{1}{\E{\max(\beta,T)}}\eqnlabel{lam}
\end{equation}
at the DP reader. We then plot the average age at the monitor as a function of $\gamma$.  Because $\E{T}=1$, the maximum update sampling rate is $\gamma=1$, which corresponds to the zero-wait policy. As $\gamma\to 0$, the average age is increasingly dominated by the average inter-read time  $1/\gamma$, because  updates become too infrequent. 

% \begin{figure}[t]
% \centering
% \includegraphics[]{MATLAB/tcnowait.eps}
% \caption{
% Average decision update age at the monitor as a function of normalized source $1$ update rate $\rho_1$ with $\rho_2 = 1 - \rho_1$ against decision update generation rate $\sig = 0.1$. The figure compares $\E{\monage}$ against a simulation.}
% \end{figure}

%RY I think your paragraph more or less collided with my explanation
%\V{since, $\E{\dpage} = \E{\maxage(t)} + \E{Z}$. We notice that if the inter-renewal time is high variance, then $\E{Z}$ is a large number. This high variance comes from variance in the inter-sample time induced by the computational machinery of the decision process. Hence, if the computational machinery of the decision process has variance, we do this trick of adding a $\beta-T$ wait to reduce the variance and get an improved $\E{Z}$, which gives us an improved $\E{\dpage}$.}



\section{Conclusion}
\label{sec:conclusion}
In this work, we focused on the problem of
timely processing of updates from multiple sources.  
Specifically, we considered a model of a publish-subscribe system where a writer publishes updates from two independent sources in a shared memory and decision updates are derived by a decision process by reading from the memory. 
% We propose a system model where the age of a decision update equals the age of oldest update used in the decision. 
The decision processing is a subscriber that works independently of how the source updates are recorded in the memory. 
%Even though the decision processing cannot improve the age in shared memory, 
Even though the decision processing operates without knowledge of the ages of updates in the shared memory,
its reading policy is  still able to improve the end-to-end decision update timeliness.

%We note that in our model, the DP Reader only has visibility into the maximum age of update pair. 
We recognize that there could be other DP reading policies that exploit knowledge of the update age processes in shared memory to further reduce decision update age at the monitor;  identifying such policies 
%the optimal policy is beyond the scope of this paper and could 
would be an interesting avenue for future research.

Another drawback of our system model is  that when  stale values are read from the shared memory, the DP will still perform its computations, even though the resulting DP update is not age-reducing. Effectively, the DP is waiting for a computation time before attempting to retrieve new updates from memory. Instead, it would be better for the DP to discard the stale updates and wait for an optimized time before reading the memory again. Analyzing such a model also remains as future work. 
% \MOVE[to conclusion]{We observe that 
% decision processing works independently of how the source updates are recorded in the memory. Hence, even though 
% decision processing cannot improve the age in shared memory, it can apply reading policies to improve the decision update timeliness.}

\section*{Acknowledgment}
This work was supported by the US National Science Foundation under grant number CNS-2148104.


% \newpage
\bibliographystyle{IEEEtran}
\bibliography{refs,AOI-2020-03}

% \newpage
\appendix


\begin{proof} (\Thmref{avg-monage-wait})
     Suppose $t\in (\tau_{i-1},\tau_i]$, the $i$th inter-sample interval. 
    We observe from Fig.~\ref{fig:tc-age-ev} that $\dpage(\tau_{i-1})=\maxage(\tau_{i-1})$ because the reader fetches update $i-1$ at that time.  However, at that time, the monitor has only received the decision update based on 
    %the computation employing 
    update $i-2$, which had age $\dpage(\tau_{i-2})$ at time $\tau_{i-2}$ and now, at time $\tau_{i-1}=\tau_{i-2}+Y_{i-1}$, has age $\dpage(\tau_{i-2})+Y_{i-1}$.  Hence, at time $\tau_{i-1}$, the monitor has age
     \begin{equation}
       \monage(\tau_{i-1})=\dpage(\tau_{i-2})+Y_{i-1}.
       \end{equation}
       Defining $H_i=\dpage(\tau_{i-2}) + Y_{i-1} - \dpage(\tau_{i-1})$, we can write       
\begin{equation}
\monage(\tau_{i-1})   =\dpage(\tau_{i-1})+ H_i.
\end{equation}
Since $\dpage(\tau_i)=\maxage(\tau_i)$ for all $i$, \begin{equation}
 H_i=\maxage(\tau_{i-2}) + Y_{i-1} - \maxage(\tau_{i-1}).
\end{equation}
It follows from  stationarity of $\maxage(t)$ and independence of the sampling times $\tau_i$  and $\maxage(t)$ that
\begin{align}
    \E{H_i}&=\E{\maxage(\tau_{i-2})} + \E{Y_{i-1}} - \E{\maxage(\tau_{i-1})}\nn
    &= \E{Y_{i-1}}=\E{T}+\E{W}.\eqnlabel{ehi} 
\end{align}
At time $\tau_{i-1}$, the $i$th busy period starts and both $\dpage(t)$ and $\monage(t)$ grow linearly at rate $1$ because neither process sees an update. Hence, $\monage(t)   =\dpage(t)+ H_i$ during the busy period. Only when the busy period completes at time $\tau_{i-1}+T_i$ does $\monage(t)$ drop and become equal to $\dpage(t)$. 
Let events $B_t$ and $I_t$ correspond to the decision process being busy and idle respectively, at time $t$.
In this interval, the event $B_t$ occurs while $\tau_{i-1} \leq t \leq \tau_{i-1} + T_i$; otherwise $I_t$ occurs if $\tau_{i-1} + T_i \leq t \leq \tau_{i}$. With these events, we can write
\begin{align}
        \monage(t) &= \begin{cases}
            \dpage(t) + H_i, & \text{if $B_t$}, \\
            \dpage(t), & \text{if $I_t$}.
        \end{cases}
    \end{align}

%RY: I rewrote this proof because the next (now commented out) statement about $h_i$ isn't correct
% where,
% $h_i = \dpage(\tau_{i-1}) - \maxage(\tau_{i-1})$. Following \eqnref{yhat-t}, $\dpage(\tau_{i-1}) = \maxage(\tau_{i-1} - Y_{i-1}) + Y_{i-1}$, we get 
% \begin{align}
%     h_i &= \maxage(\tau_{i-1} - Y_{i-1}) + Y_{i-1} -\maxage(\tau_{i-1}) \nonumber \\
%         &= \maxage(\tau_{i-2}) + Y_{i-1} - \maxage(\tau_{i-1}).
% \end{align}
% With $Y_{i-1} \equiv Y$, and from the stationarity of $\maxage(t)$ age process,
% \begin{equation}
%     \E{h_i} = \E{\maxage(\tau_{i-2})} + \E{Y_{i-1}} - \E{\maxage(\tau_{i-1})} = \E{Y} \eqnlabel{dp-height}. 
% \end{equation}
For $t\ge \tau_{i-1}$, event $B_t$ is independent of $H_i$, and it follows from  the law of total expectation that
% \RY{Let's rename $A_i$ as $B$ (corresponding to the DP is ``Busy'' and $B_i$ as $I$, corresponding to the DP is ``Idle''    I turned on subequations for the moment to say (24a) is wrong. When you see the bug there, you will realize that you can simplify the steps A LOT. In fact, you never need to mention $\xhat$. By the way, you don't need commas on each line because it's a single logical chain, i.e. you wouldn't write a=b=c=d=... as a=b,=c,=d,=...} 
  \begin{align}
        \E{\monage(t)} &= \E{\dpage(t) +H _i| B}\prob{B_t} + \E{\dpage(t) | I_t} \prob{I_t} \nonumber \\
 %       &= \E{\dpage(t) | B_t}P(B_t) + \E{h_i|B_t}\prob{B_t} + \E{\dpage(t)|I_t}\prob{I_t} \nonumber \\     
        &= \E{\dpage(t)} + \E{H_i}\prob{B_t}.
        % \nonumber \\
        % &= \E{\maxage(t - Z(t)) + Z(t)} + \E{h_i}\prob{B} \nonumber \\
        % &= \E{\maxage(t - Z(t))} + \E{Z} + \E{h_i}\prob{B}. 
        \eqnlabel{ezhat}
    \end{align}
    % With $T_i  \equiv T, W_i  \equiv W$, $Y_i = T_i + W_i \equiv T+W $, and since $T$ and $W$ are independent, from \eqnref{bcwexc-mom},
    % \begin{equation}
    %     \E{Z} = \frac{1}{2}(\E{T} + \E{W} + \frac{\Var{T} + \Var{W}}{\E{T} + \E{W}}). \eqnlabel{ez}
    % \end{equation}
    In each renewal period, the decision process is busy for time $T$ and then idle for time $W$. By considering a renewal reward process in which a reward  $T$ is earned for the busy period, it follows that the limiting fraction of time spent in a busy state is given by
    \begin{equation}
        \prob{B_t} = \frac{\E{T}}{\E{T} + \E{W}}. \eqnlabel{pb}
    \end{equation}
    Applying \eqnref{ehi} and \eqnref{pb} to \eqnref{ezhat} yields the claim.
    % \begin{align}
    %     \E{\monage(t)} &= \E{\dpage(t)} + \E{T}
    % \end{align}
    % Using \eqnref{avg-dpage} in above equation, \eqnref{tmp1} follows, and observing $\E{Y^2} = \Var{Y} + (\E{Y})^2$, \eqnref{avg-monage-wait} follows.
    % Solving \eqnref{ezhat} using \eqnref{ez}, \eqnref{ehi}, and \eqnref{pb}, \eqnref{avg-monage-wait} follows.
\end{proof}

\end{document}

